\section{Main}
\label{sec:main}

Human “genetic support” implicating genes in disease substantially increases drug target viability and provides a strong justification for experimental studies [Plenge, Minikel, Nelson]. Large-scale genome-wide association studies (GWAS) and exome-wide association studies (ExWAS) provide a strong foundation for this effort. These studies identify hundreds of thousands of variants associated with human traits. However, researchers struggle to determine the specific genes implicated by these associations. Common variant associations often fall in non-coding regions with unclear links to effector genes. Conversely, rare variant associations from ExWAS have limited statistical power. Furthermore, association studies typically report only highly significant associations. These stringent thresholds offer little information on the extent to which sub-significant associations provide genomic support for a gene and trait of interest.

Work to address these challenges mainly falls into three areas. First, effector gene prediction methods attempt to localize significant genome-wide association study (GWAS) signals to specific genes. Fine-mapping methods such as SuSiE[cite] localize GWAS signals to causal variants using sparse priors to calculate better posterior distributions. Fine-mapping is traditionally based on properties of variant association strength but can also be improved by leveraging functional annotations, including ATAC-seq[Corces, Bi], as demonstrated by tools like PolyFun[cite] and fgwas[cite]. These variants can then be linked to putative effector genes using variant-to-gene (V2G) scores that integrate multiple types of evidence (e.g. Hi-C[Jung], cS2G[Gazal]). These predictions are often manually combined into “effector gene lists” that accompany GWAS papers, although these lists lack standardization and rigor[Costanzo]. Second, gene-level approaches such as MAGMA[cite] and MetaXcan[cite] assign genetic support statistics to all genes in the genome by integrating multiple SNPs with (potentially) eQTLs. Third, data from ExWAS, in which burden[cite] or variance component tests[cite] aggregate heuristic masks[Trang] of variants can directly implicate disease-relevant genes, albeit with significantly less power than GWAS[19, Dornbos]. These data and methods each address various parts of the genetic support question, but their heterogeneous assumptions and mathematical frameworks make their integration difficult[HuGE].

To navigate this complexity, we developed a Bayesian statistical method named FALCON (Framework for Analysis of Linkage and Combined Omics using Networks). This model integrates association studies—taking into account correlation— with functional genomic datasets (functional annotations and V2G) to calculate genetic support. We built the network's architecture upon established genetic algorithms (e.g. SuSiE[cite]) and used empirically generated data as priors to inform this structure (e.g. LD Score regression[cite]). We validated the accuracy of this framework on both simulations and real data, using these results to identify the genetic and genomic features most useful for predicting genetic support. We applied FALCON to 1370 full GWAS summary statistics with a total of 117,748 curated GWAS lead SNPs to produce a Genetic Map (G-Map). We used this map to (a) conduct a retrospective analysis of drug target relative success; (b) uncover distinct tissue-mediated mechanisms between SNPs and genes; (c) and evaluate the utility of GWAS associations for prioritizing rare disease genes. This method and map provides a foundation to accurately estimate probabilities of genetic support between every gene and every complex trait.


